\begin{pastexam}[label={px:ncku-112-prog-15}]{成大資訊聯招 112 學年度 第 15 題}
（計算）

(10\%) The following graph is about the maximum flow problem. Each edge is labeled with its capacity.
What is the maximum flow from vertex $s$ to vertex $t$.

\begin{center}
\begin{tikzpicture}[x=1cm,y=1cm,>={Stealth[length=2.2mm]},
  v/.style={circle,draw,thick,inner sep=0pt,minimum size=8mm,font=\sffamily\small},
  w/.style={font=\sffamily\small,inner sep=2pt}]
\node[v] (s) at (0,0) {\textit{s}};
\node[v] (n1) at (2.5,1.5) {1};
\node[v] (n3) at (5,1.5) {3};
\node[v] (n2) at (2.5,-1.5) {2};
\node[v] (n4) at (5,-1.5) {4};
\node[v] (t) at (7.5,0) {\textit{t}};
\draw[->] (s)--node[w,above left]{18}(n1);
\draw[->] (s)--node[w,below left]{16}(n2);
\draw[->] (n2)--node[w,left,pos=0.45]{6}(n1);
\draw[->] (n1)--node[w,above]{16}(n3);
\draw[->] (n3)--node[w,above left,pos=0.35]{19}(n2);
\draw[->] (n2)--node[w,below]{14}(n4);
\draw[->] (n4)--node[w,right,pos=0.45]{7}(n3);
\draw[->] (n3)--node[w,above right]{22}(t);
\draw[->] (n4)--node[w,below right]{8}(t);
\end{tikzpicture}
\end{center}
\end{pastexam}

\begin{pastsoln}
答案：最大流是 $30$。 下面用 Edmonds--Karp 求出一個值 $30$ 的流，再用一個容量 $30$ 的割證明
它是最大的（資結書第~26.4、26.5~節）。 BFS 依 $s, 1, 2, 3, 4, t$ 的順序檢查鄰
居。 這張網路沒有反平行邊（$1 \to 3$、$3 \to 2$、$2 \to 1$ 是一個方向一致的環，不是反平行邊），
可以直接用。

\begin{center}
\small
\begin{tabular}{c|cccccc|lcc}
\toprule
輪 & $s$ & $1$ & $2$ & $3$ & $4$ & $t$ & 增廣路徑 & 瓶頸 & $|f|$ \\
\midrule
1 & 0 & 1 & 1 & 2 & 2 & 3 & $s \to 1 \to 3 \to t$ & $16$ & $16$ \\
2 & 0 & 1 & 1 & 3 & 2 & 3 & $s \to 2 \to 4 \to t$ & $8$ & $24$ \\
3 & 0 & 1 & 1 & 3 & 2 & 4 & $s \to 2 \to 4 \to 3 \to t$ & $6$ & $30$ \\
4 & 0 & 1 & 1 & $\infty$ & $\infty$ & $\infty$ & 走不到 $t$ & & \\
\bottomrule
\end{tabular}
\end{center}

每一輪的細節：
\begin{itemize}
\item \textbf{第 1 輪。} 剩餘網路就是原圖。 $s$ 發現 $1$、$2$；$1$ 發現 $3$；$2$ 發現 $4$
（$2 \to 1$ 指向已發現的 $1$）；$3$ 發現 $t$。 路徑 $s \to 1 \to 3 \to t$，剩餘容量 $18, 16, 22$，
瓶頸 $16$。 增廣後 $1 \to 3$ 滿了，多出反向邊 $1 \to s$（$16$）、$3 \to 1$（$16$）、$t \to 3$
（$16$）。
\item \textbf{第 2 輪。} $1$ 的剩餘邊只剩反向邊 $1 \to s$，走不到 $3$。 $2$ 發現 $4$（距離
$2$），$4$ 發現 $3$ 和 $t$（距離 $3$）。 路徑 $s \to 2 \to 4 \to t$，剩餘容量 $16, 14, 8$，瓶頸
$8$。 增廣後 $4 \to t$ 滿了。
\item \textbf{第 3 輪。} $4$ 到不了 $t$，只能經 $4 \to 3$（$7$）到 $3$，再走 $3 \to t$（剩
$22 - 16 = 6$）。 路徑 $s \to 2 \to 4 \to 3 \to t$，剩餘容量 $16 - 8 = 8$、$14 - 8 = 6$、$7$、$6$，
瓶頸 $6$。 增廣後 $2 \to 4$、$3 \to t$ 都滿了。
\item \textbf{第 4 輪。} $s$ 出去還有 $s \to 1$（$2$）、$s \to 2$（$2$）。 $1$ 的出邊 $1 \to 3$
滿了；$2$ 的出邊 $2 \to 1$ 指向已發現的 $1$，$2 \to 4$ 滿了；$3 \to 2$、$2 \to 1$ 上沒有流，
沒有反向邊可走。 BFS 只走到 $s, 1, 2$，找不到增廣路徑。
\end{itemize}

最後的流（流量$/$容量）：$s \to 1$：$16/18$、$s \to 2$：$14/16$、$2 \to 1$：$0/6$、$1 \to 3$：
$16/16$、$3 \to 2$：$0/19$、$2 \to 4$：$14/14$、$4 \to 3$：$6/7$、$3 \to t$：$22/22$、$4 \to t$：
$8/8$。 逐點檢查守恆：$1$ 進 $16$ 出 $16$；$2$ 進 $14$ 出 $14$；$3$ 進 $16 + 6$ 出 $22$；$4$ 進
$14$ 出 $6 + 8$。 值 $16 + 14 = 30$，等於流進 $t$ 的 $22 + 8$。

\textbf{驗證是最大流。} 最後一輪 BFS 走得到的是 $S = \{s, 1, 2\}$，$T = \{3, 4, t\}$。 從 $S$
指向 $T$ 的邊是 $1 \to 3$（$16$）和 $2 \to 4$（$14$），都滿了；從 $T$ 指回 $S$ 的邊只有
$3 \to 2$，流量 $0$。 $c(S, T) = 16 + 14 = 30 = |f|$，由資結書系理~26.5，沒有任何流能超過
$30$。
\end{pastsoln}
